\documentclass[11pt,a4paper]{article}
\usepackage[utf8]{inputenc}
\usepackage[T1]{fontenc}
\usepackage[catalan]{babel}
\usepackage{lmodern}
\usepackage[sfdefault,lining]{FiraSans}
\usepackage[a4paper,top=2.3cm,bottom=2.2cm,left=1.8cm,right=1.8cm,headheight=1.3cm,headsep=0.4cm,footskip=1cm]{geometry}
\usepackage{amsmath,amssymb}
\usepackage{icomma}
\usepackage{graphicx}
\usepackage[table]{xcolor}
\usepackage{tikz}
\usetikzlibrary{calc}
\usepackage{fontawesome5}
\usepackage{ccicons}
\usepackage{fancyhdr}
\usepackage{lastpage}
\usepackage{tasks}
\usepackage{needspace}
\usepackage[most]{tcolorbox}
\usepackage{array}
\usepackage[hidelinks,pdftitle={Fitxa de repàs: Global de Programació lineal. 2n de Batxillerat CCSS},pdfsubject={Matemàtiques Aplicades a les Ciències Socials II},pdfkeywords={programació lineal, regió factible, vèrtexs, infinites solucions, regions no fitades, PAU, repàs}]{hyperref}

\definecolor{navy}{HTML}{1B2748}
\definecolor{accent}{HTML}{E8603A}
\definecolor{teal}{HTML}{1F9E90}
\definecolor{soft}{HTML}{EEF1F8}
\definecolor{muted}{HTML}{5D6685}

\setlength{\parindent}{0pt}
\setlength{\parskip}{3pt}
\settasks{label=\alph*), label-format=\bfseries\color{navy}, label-width=1.3em, item-indent=1.8em, column-sep=0.6em, after-item-skip=5pt, before-skip=3pt, after-skip=4pt}

% ---------- capçalera i peu ----------
\newcommand{\llicencia}{\href{https://creativecommons.org/licenses/by-nc-sa/4.0/deed.ca}{\ccbyncsa}\ \ Llicència Creative Commons Reconeixement-NoComercial-CompartirIgual 4.0 Internacional (CC BY-NC-SA 4.0)}
\pagestyle{fancy}
\fancyhf{}
\fancyhead[L]{\raisebox{-0.35\height}{\includegraphics[height=0.95cm]{logo-gva.png}}}
\fancyhead[C]{\small\color{muted}Fitxa de repàs · Global de Programació lineal · 2n Batxillerat CCSS}
\fancyhead[R]{\raisebox{-0.35\height}{\includegraphics[height=0.95cm]{logo-vermellar.png}}}
\fancyfoot[L]{\scriptsize\color{muted}\llicencia}
\fancyfoot[R]{\small\color{muted}\thepage}
\renewcommand{\headrulewidth}{0.4pt}
\renewcommand{\headrule}{\color{muted!40}\hrule height \headrulewidth}
\fancypagestyle{primera}{\fancyhf{}\renewcommand{\headrulewidth}{0pt}
  \fancyfoot[L]{\scriptsize\color{muted}\llicencia}
  \fancyfoot[R]{\small\color{muted}\thepage}}

% ---------- nivells i exercicis ----------
\newcommand{\estrelles}[1]{{\color{accent}\ifcase#1\or\faStar\faStar[regular]\faStar[regular]\or\faStar\faStar\faStar[regular]\or\faStar\faStar\faStar\fi}}
\newcounter{exer}
\newcommand{\ex}[1]{\par\needspace{4\baselineskip}\medskip\stepcounter{exer}\noindent{\color{navy}\bfseries\theexer.}\ \ }

\newcommand{\seccio}[2]{\par\needspace{10\baselineskip}\bigskip
  \begin{tcolorbox}[colback=navy,colframe=navy,coltext=white,arc=2mm,boxrule=0pt,left=3mm,right=3mm,top=1.2mm,bottom=1.2mm]
  {\large\bfseries #1\quad #2}
  \end{tcolorbox}\vspace{-2pt}}


\newcommand{\recorda}[1]{\par\smallskip\begin{tcolorbox}[colback=accent!7,colframe=accent!55,arc=2mm,boxrule=0.5pt,left=2.5mm,right=2.5mm,top=1mm,bottom=1mm]\small{\color{accent}\faLightbulb[regular]}\ \textbf{Recorda:}\ #1\end{tcolorbox}}

% =====================================================================
% Capçalera FentMates de les fitxes: un color per a cada curs.
% Els colors són els de la pàgina inicial (LEVELS a assets/js/materials.js);
% el logo ve de _eines/logo/pdf/ (python3 _eines/logo/genera.py).
%
% Ús, al final del preàmbul de la unitat:
%   % =====================================================================
% Capçalera FentMates de les fitxes: un color per a cada curs.
% Els colors són els de la pàgina inicial (LEVELS a assets/js/materials.js);
% el logo ve de _eines/logo/pdf/ (python3 _eines/logo/genera.py).
%
% Ús, al final del preàmbul de la unitat:
%   % =====================================================================
% Capçalera FentMates de les fitxes: un color per a cada curs.
% Els colors són els de la pàgina inicial (LEVELS a assets/js/materials.js);
% el logo ve de _eines/logo/pdf/ (python3 _eines/logo/genera.py).
%
% Ús, al final del preàmbul de la unitat:
%   \input{../comu/fentmates.tex}
%   \fmcurs{1eso}
% i, al començament del document:
%   \capcalera{Matemàtiques}{1r ESO}{Fitxa d'exercicis · Unitat 2}{Divisibilitat}
% Opcional abans de \input: \def\fmneed{14} (espai mínim abans de cada secció).
% =====================================================================
\makeatletter
\def\fm@logodir{../../logo/pdf}
\providecommand{\fmneed}{10}
\def\fm@cursnom{}

% clau del curs, color i color del text damunt del color
\newcommand{\fm@curs}[3]{\@namedef{fm@color@#1}{#2}\@namedef{fm@text@#1}{#3}}
\fm@curs{1eso}{E6B84A}{navy}
\fm@curs{2eso}{3B82F6}{white}
\fm@curs{3eso}{FF7A50}{white}
\fm@curs{4eso-a}{A78BFA}{white}
\fm@curs{4eso-b}{6D28D9}{white}
\fm@curs{1batx-mi}{9B1C1C}{white}
\fm@curs{1batx-ccss}{EF6B6B}{white}
\fm@curs{2batx-mii}{1E6B47}{white}
\fm@curs{2batx-ccss}{2BB3A3}{white}

\newcommand{\fmcurs}[1]{%
  \@ifundefined{fm@color@#1}{\PackageError{fentmates}{Curs desconegut: #1}{}}{}%
  \edef\fm@tmp{\noexpand\definecolor{curs}{HTML}{\@nameuse{fm@color@#1}}%
               \noexpand\colorlet{curstext}{\@nameuse{fm@text@#1}}}\fm@tmp
  \colorlet{cursfosc}{curs!72!black}%
  \def\fm@logo{\fm@logodir/vertical-#1.pdf}%
  \def\fm@marca{\fm@logodir/marca-#1.pdf}}

% ---------- capçalera de la primera pàgina ----------
% #1 matèria (petit), #2 curs, #3 tipus de document i unitat, #4 títol
\newcommand{\capcalera}[4]{%
  \gdef\fm@cursnom{#2}%
  \par\noindent
  \begin{tcolorbox}[colback=curs,colframe=curs,boxrule=0pt,arc=2.2mm,width=3.5cm,height=2.15cm,
      valign=center,left=3mm,right=2mm,top=1mm,bottom=1mm,nobeforeafter,box align=center]
    \color{curstext}{\footnotesize #1}\par\vspace{1pt}{\fontsize{19}{21}\selectfont\bfseries #2\par}%
  \end{tcolorbox}\hspace{5mm}%
  \begin{minipage}[c]{\dimexpr\linewidth-3.5cm-5mm-2.9cm\relax}
    {\color{cursfosc}\small\bfseries\MakeUppercase{#3}}\par\vspace{2pt}
    {\color{navy}\fontsize{20}{23}\selectfont\bfseries #4\par}
  \end{minipage}\hfill
  \begin{minipage}[c]{2.6cm}\raggedleft\includegraphics[width=2.6cm]{\fm@logo}\end{minipage}\par
  \vspace{2pt}}

% ---------- barres de secció del color del curs ----------
\renewcommand{\seccio}[2]{\par\needspace{\fmneed\baselineskip}\bigskip
  \begin{tcolorbox}[colback=curs,colframe=curs,coltext=curstext,arc=2mm,boxrule=0pt,left=3mm,right=3mm,top=1.2mm,bottom=1.2mm]
  {\large\bfseries #1\quad #2}
  \end{tcolorbox}\vspace{-2pt}}

% ---------- estrelles de dificultat del color del curs (a Batxillerat no se'n posen) ----------
\@ifundefined{estrelles}{}{%
  \renewcommand{\estrelles}[1]{{\color{curs}\ifcase#1\or\faStar\faStar[regular]\faStar[regular]\or\faStar\faStar\faStar[regular]\or\faStar\faStar\faStar\fi}}}

% ---------- capçalera i peu de les pàgines ----------
% primera pàgina: logos institucionals al peu
\fancypagestyle{primera}{\fancyhf{}\renewcommand{\headrulewidth}{0pt}
  \fancyfoot[L]{\raisebox{-0.32\height}{\includegraphics[height=0.78cm]{logo-gva.png}}\hspace{5mm}%
    \raisebox{-0.32\height}{\includegraphics[height=0.78cm]{logo-vermellar.png}}}
  \fancyfoot[R]{\scriptsize\color{muted}\href{https://creativecommons.org/licenses/by-nc-sa/4.0/deed.ca}{\ccbyncsa}\ \ CC BY-NC-SA 4.0 · fentmates.org\qquad\small\thepage}}
% la resta: el símbol a l'esquerra i el curs a la dreta, amb el títol al centre com abans
\fancyhead[L]{\raisebox{-0.3\height}{\includegraphics[height=0.62cm]{\fm@marca}}}
\fancyhead[R]{\small\bfseries\color{cursfosc}\fm@cursnom}
\renewcommand{\headrule}{\color{curs!45}\hrule height \headrulewidth}
\makeatother

%   \fmcurs{1eso}
% i, al començament del document:
%   \capcalera{Matemàtiques}{1r ESO}{Fitxa d'exercicis · Unitat 2}{Divisibilitat}
% Opcional abans de \input: \def\fmneed{14} (espai mínim abans de cada secció).
% =====================================================================
\makeatletter
\def\fm@logodir{../../logo/pdf}
\providecommand{\fmneed}{10}
\def\fm@cursnom{}

% clau del curs, color i color del text damunt del color
\newcommand{\fm@curs}[3]{\@namedef{fm@color@#1}{#2}\@namedef{fm@text@#1}{#3}}
\fm@curs{1eso}{E6B84A}{navy}
\fm@curs{2eso}{3B82F6}{white}
\fm@curs{3eso}{FF7A50}{white}
\fm@curs{4eso-a}{A78BFA}{white}
\fm@curs{4eso-b}{6D28D9}{white}
\fm@curs{1batx-mi}{9B1C1C}{white}
\fm@curs{1batx-ccss}{EF6B6B}{white}
\fm@curs{2batx-mii}{1E6B47}{white}
\fm@curs{2batx-ccss}{2BB3A3}{white}

\newcommand{\fmcurs}[1]{%
  \@ifundefined{fm@color@#1}{\PackageError{fentmates}{Curs desconegut: #1}{}}{}%
  \edef\fm@tmp{\noexpand\definecolor{curs}{HTML}{\@nameuse{fm@color@#1}}%
               \noexpand\colorlet{curstext}{\@nameuse{fm@text@#1}}}\fm@tmp
  \colorlet{cursfosc}{curs!72!black}%
  \def\fm@logo{\fm@logodir/vertical-#1.pdf}%
  \def\fm@marca{\fm@logodir/marca-#1.pdf}}

% ---------- capçalera de la primera pàgina ----------
% #1 matèria (petit), #2 curs, #3 tipus de document i unitat, #4 títol
\newcommand{\capcalera}[4]{%
  \gdef\fm@cursnom{#2}%
  \par\noindent
  \begin{tcolorbox}[colback=curs,colframe=curs,boxrule=0pt,arc=2.2mm,width=3.5cm,height=2.15cm,
      valign=center,left=3mm,right=2mm,top=1mm,bottom=1mm,nobeforeafter,box align=center]
    \color{curstext}{\footnotesize #1}\par\vspace{1pt}{\fontsize{19}{21}\selectfont\bfseries #2\par}%
  \end{tcolorbox}\hspace{5mm}%
  \begin{minipage}[c]{\dimexpr\linewidth-3.5cm-5mm-2.9cm\relax}
    {\color{cursfosc}\small\bfseries\MakeUppercase{#3}}\par\vspace{2pt}
    {\color{navy}\fontsize{20}{23}\selectfont\bfseries #4\par}
  \end{minipage}\hfill
  \begin{minipage}[c]{2.6cm}\raggedleft\includegraphics[width=2.6cm]{\fm@logo}\end{minipage}\par
  \vspace{2pt}}

% ---------- barres de secció del color del curs ----------
\renewcommand{\seccio}[2]{\par\needspace{\fmneed\baselineskip}\bigskip
  \begin{tcolorbox}[colback=curs,colframe=curs,coltext=curstext,arc=2mm,boxrule=0pt,left=3mm,right=3mm,top=1.2mm,bottom=1.2mm]
  {\large\bfseries #1\quad #2}
  \end{tcolorbox}\vspace{-2pt}}

% ---------- estrelles de dificultat del color del curs (a Batxillerat no se'n posen) ----------
\@ifundefined{estrelles}{}{%
  \renewcommand{\estrelles}[1]{{\color{curs}\ifcase#1\or\faStar\faStar[regular]\faStar[regular]\or\faStar\faStar\faStar[regular]\or\faStar\faStar\faStar\fi}}}

% ---------- capçalera i peu de les pàgines ----------
% primera pàgina: logos institucionals al peu
\fancypagestyle{primera}{\fancyhf{}\renewcommand{\headrulewidth}{0pt}
  \fancyfoot[L]{\raisebox{-0.32\height}{\includegraphics[height=0.78cm]{logo-gva.png}}\hspace{5mm}%
    \raisebox{-0.32\height}{\includegraphics[height=0.78cm]{logo-vermellar.png}}}
  \fancyfoot[R]{\scriptsize\color{muted}\href{https://creativecommons.org/licenses/by-nc-sa/4.0/deed.ca}{\ccbyncsa}\ \ CC BY-NC-SA 4.0 · fentmates.org\qquad\small\thepage}}
% la resta: el símbol a l'esquerra i el curs a la dreta, amb el títol al centre com abans
\fancyhead[L]{\raisebox{-0.3\height}{\includegraphics[height=0.62cm]{\fm@marca}}}
\fancyhead[R]{\small\bfseries\color{cursfosc}\fm@cursnom}
\renewcommand{\headrule}{\color{curs!45}\hrule height \headrulewidth}
\makeatother

%   \fmcurs{1eso}
% i, al començament del document:
%   \capcalera{Matemàtiques}{1r ESO}{Fitxa d'exercicis · Unitat 2}{Divisibilitat}
% Opcional abans de \input: \def\fmneed{14} (espai mínim abans de cada secció).
% =====================================================================
\makeatletter
\def\fm@logodir{../../logo/pdf}
\providecommand{\fmneed}{10}
\def\fm@cursnom{}

% clau del curs, color i color del text damunt del color
\newcommand{\fm@curs}[3]{\@namedef{fm@color@#1}{#2}\@namedef{fm@text@#1}{#3}}
\fm@curs{1eso}{E6B84A}{navy}
\fm@curs{2eso}{3B82F6}{white}
\fm@curs{3eso}{FF7A50}{white}
\fm@curs{4eso-a}{A78BFA}{white}
\fm@curs{4eso-b}{6D28D9}{white}
\fm@curs{1batx-mi}{9B1C1C}{white}
\fm@curs{1batx-ccss}{EF6B6B}{white}
\fm@curs{2batx-mii}{1E6B47}{white}
\fm@curs{2batx-ccss}{2BB3A3}{white}

\newcommand{\fmcurs}[1]{%
  \@ifundefined{fm@color@#1}{\PackageError{fentmates}{Curs desconegut: #1}{}}{}%
  \edef\fm@tmp{\noexpand\definecolor{curs}{HTML}{\@nameuse{fm@color@#1}}%
               \noexpand\colorlet{curstext}{\@nameuse{fm@text@#1}}}\fm@tmp
  \colorlet{cursfosc}{curs!72!black}%
  \def\fm@logo{\fm@logodir/vertical-#1.pdf}%
  \def\fm@marca{\fm@logodir/marca-#1.pdf}}

% ---------- capçalera de la primera pàgina ----------
% #1 matèria (petit), #2 curs, #3 tipus de document i unitat, #4 títol
\newcommand{\capcalera}[4]{%
  \gdef\fm@cursnom{#2}%
  \par\noindent
  \begin{tcolorbox}[colback=curs,colframe=curs,boxrule=0pt,arc=2.2mm,width=3.5cm,height=2.15cm,
      valign=center,left=3mm,right=2mm,top=1mm,bottom=1mm,nobeforeafter,box align=center]
    \color{curstext}{\footnotesize #1}\par\vspace{1pt}{\fontsize{19}{21}\selectfont\bfseries #2\par}%
  \end{tcolorbox}\hspace{5mm}%
  \begin{minipage}[c]{\dimexpr\linewidth-3.5cm-5mm-2.9cm\relax}
    {\color{cursfosc}\small\bfseries\MakeUppercase{#3}}\par\vspace{2pt}
    {\color{navy}\fontsize{20}{23}\selectfont\bfseries #4\par}
  \end{minipage}\hfill
  \begin{minipage}[c]{2.6cm}\raggedleft\includegraphics[width=2.6cm]{\fm@logo}\end{minipage}\par
  \vspace{2pt}}

% ---------- barres de secció del color del curs ----------
\renewcommand{\seccio}[2]{\par\needspace{\fmneed\baselineskip}\bigskip
  \begin{tcolorbox}[colback=curs,colframe=curs,coltext=curstext,arc=2mm,boxrule=0pt,left=3mm,right=3mm,top=1.2mm,bottom=1.2mm]
  {\large\bfseries #1\quad #2}
  \end{tcolorbox}\vspace{-2pt}}

% ---------- estrelles de dificultat del color del curs (a Batxillerat no se'n posen) ----------
\@ifundefined{estrelles}{}{%
  \renewcommand{\estrelles}[1]{{\color{curs}\ifcase#1\or\faStar\faStar[regular]\faStar[regular]\or\faStar\faStar\faStar[regular]\or\faStar\faStar\faStar\fi}}}

% ---------- capçalera i peu de les pàgines ----------
% primera pàgina: logos institucionals al peu
\fancypagestyle{primera}{\fancyhf{}\renewcommand{\headrulewidth}{0pt}
  \fancyfoot[L]{\raisebox{-0.32\height}{\includegraphics[height=0.78cm]{logo-gva.png}}\hspace{5mm}%
    \raisebox{-0.32\height}{\includegraphics[height=0.78cm]{logo-vermellar.png}}}
  \fancyfoot[R]{\scriptsize\color{muted}\href{https://creativecommons.org/licenses/by-nc-sa/4.0/deed.ca}{\ccbyncsa}\ \ CC BY-NC-SA 4.0 · fentmates.org\qquad\small\thepage}}
% la resta: el símbol a l'esquerra i el curs a la dreta, amb el títol al centre com abans
\fancyhead[L]{\raisebox{-0.3\height}{\includegraphics[height=0.62cm]{\fm@marca}}}
\fancyhead[R]{\small\bfseries\color{cursfosc}\fm@cursnom}
\renewcommand{\headrule}{\color{curs!45}\hrule height \headrulewidth}
\makeatother

\fmcurs{2batx-ccss}
\newcommand{\res}[1]{\boxed{#1}}\newcommand{\rest}[1]{\fbox{#1}}
\begin{document}
\thispagestyle{primera}

\capcalera{Matemàtiques CCSS}{2n Batx}{Fitxa de repàs · Global}{Programació lineal}

\vspace{10pt}
Nom i cognoms: \rule{8.4cm}{0.4pt}\hfill Grup: \rule{1.6cm}{0.4pt}\hfill Data: \rule{2.2cm}{0.4pt}

\vspace{8pt}
\begin{tcolorbox}[colback=soft,colframe=soft,arc=2mm,left=3mm,right=3mm,top=1.5mm,bottom=1.5mm]
\small Quatre problemes de repàs de tota la unitat 3, amb el nivell de la PAU: regions factibles i vèrtexs, optimització en regions fitades i no fitades, infinites solucions, solucions enteres, plantejament de problemes i lectura d'una regió a partir d'un gràfic. Fes les gràfiques amb regle i paper quadriculat, justifica cada pas i, quan acabes, comprova-ho amb les solucions del final.
\end{tcolorbox}

\seccio{1}{Regió factible, optimització i infinites solucions}

\ex{1} Es considera la regió del pla determinada pel sistema d'inequacions\[-x+y\leq 3,\qquad x+2y\leq 12,\qquad 2x-y\leq 9,\qquad y\geq -1.\]\begin{tasks}(1)
  \task Representa gràficament la regió i calcula'n els vèrtexs.
  \task Justifica si els punts $P(0,4)$ i $Q(3,2)$ pertanyen a la regió.
  \task Determina el valor màxim i el valor mínim de $f(x,y)=3x+2y$ en la regió i els punts on s'assoleixen.
  \task Es considera la funció $g(x,y)=kx+2y$, amb $k\in\mathbb{R}$. Troba els valors de $k$ perquè el \textbf{màxim} de $g$ en la regió s'assolisca en tots els punts d'un costat. Per a cada valor, indica el costat i el valor màxim.
\end{tasks}\recorda{Si les rectes de nivell de la funció objectiu són paral·leles a un costat de la regió, l'òptim pot assolir-se en tots els punts d'aquest costat: hi ha infinites solucions.}


\seccio{2}{Problema de maximitzar amb solucions enteres}

\ex{1} Una cooperativa citrícola prepara dos tipus de caixes de regal. Cada caixa del tipus A porta 2~kg de taronges i 1~kg de mandarines, i cada caixa del tipus B porta 1~kg de taronges i 2~kg de mandarines. La cooperativa disposa de 100~kg de taronges i de 80~kg de mandarines. Per a atendre les comandes ha de preparar, com a mínim, 20 caixes en total, i almenys 10 han de ser del tipus B. Per cada caixa del tipus A guanya 5~€ i per cada caixa del tipus B, 4~€.\begin{tasks}(1)
  \task Planteja el problema de programació lineal que permet maximitzar el benefici: defineix les variables, la funció objectiu i les restriccions.
  \task Representa la regió factible i calcula'n els vèrtexs.
  \task Quantes caixes de cada tipus ha de preparar per a obtindre el benefici màxim? Quin és aquest benefici?
  \task Si el benefici passara a ser de 6~€ per caixa A i de 3~€ per caixa B, de quantes maneres es podria obtindre el benefici màxim? Indica-les totes i el benefici.
\end{tasks}\recorda{Si les variables compten objectes (caixes, persones…), les solucions han de ser nombres enters: quan l'òptim és tot un segment, busca'n els punts amb coordenades enteres.}


\seccio{3}{Problema de minimitzar en una regió no fitada}

\ex{1} Un ramader alimenta el seu bestiar amb dos pinsos, P1 i P2. Cada sac de P1 aporta 3 unitats de proteïnes, 1 de fibra i 1 de minerals; cada sac de P2 aporta 1 unitat de proteïnes, 1 de fibra i 3 de minerals. Cada setmana el bestiar necessita, com a mínim, 15 unitats de proteïnes, 9 de fibra i 15 de minerals. Un sac de P1 costa 40~€ i un sac de P2, 30~€.\begin{tasks}(1)
  \task Planteja el problema que permet trobar quants sacs de cada pinso ha de comprar setmanalment per a cobrir les necessitats amb el cost mínim.
  \task Representa la regió factible i calcula'n els vèrtexs. És una regió fitada?
  \task Resol el problema: quants sacs de cada tipus ha de comprar i quin és el cost mínim?
  \task Té sentit preguntar-se pel cost màxim? Raona-ho. Si el ramader comprara 4 sacs de cada pinso, cobriria les necessitats?
\end{tasks}\recorda{En una regió no fitada l'òptim pot no existir. Si els coeficients de la funció objectiu són positius i la regió s'estén cap amunt i cap a la dreta, hi ha mínim però no màxim.}


\seccio{4}{Regió donada en un gràfic}

\ex{1} La regió ombrejada és la regió factible d'un problema de programació lineal. Els vèrtexs tenen coordenades enteres.\begin{center}\begin{tikzpicture}[x=0.4600cm,y=0.4500cm,baseline=(current bounding box.center)]\begin{scope}\clip (-1,-1) rectangle (9,7);\foreach \i in {0,2,...,9} \draw[muted!25,very thin] (\i,-1)--(\i,7);\foreach \i in {0,2,...,7} \draw[muted!25,very thin] (-1,\i)--(9,\i);\fill[teal!22] (7.0000,3.0000)--(4.0000,6.0000)--(0.0000,6.0000)--(0.0000,2.0000)--(4.0000,0.0000)--cycle;\end{scope}\draw[-latex,thick] (-1,0)--(9.5,0) node[below left=-1pt,font=\scriptsize]{$x$};\draw[-latex,thick] (0,-1)--(0,7.5) node[below left=-1pt,font=\scriptsize]{$y$};\node[below,font=\tiny,text=muted] at (2,0) {$2$};\node[below,font=\tiny,text=muted] at (4,0) {$4$};\node[below,font=\tiny,text=muted] at (6,0) {$6$};\node[below,font=\tiny,text=muted] at (8,0) {$8$};\node[left,font=\tiny,text=muted] at (0,2) {$2$};\node[left,font=\tiny,text=muted] at (0,4) {$4$};\node[left,font=\tiny,text=muted] at (0,6) {$6$};\begin{scope}\clip (-1,-1) rectangle (9,7);\draw[accent,line width=1.1pt] (-2.0000,3.0000)--(10.0000,-3.0000);\draw[blue!70!black,line width=1.1pt] (-2.0000,6.0000)--(10.0000,6.0000);\draw[violet,line width=1.1pt] (12.0000,-2.0000)--(2.0000,8.0000);\draw[orange!85!black,line width=1.1pt] (2.0000,-2.0000)--(12.0000,8.0000);\end{scope}\filldraw[fill=white,draw=navy] (0.000,2.000) circle (1.6pt);\node[font=\scriptsize\bfseries,text=navy] at ($(0.000,2.000)+(-155.0:0.32cm)$) {A};\filldraw[fill=white,draw=navy] (4.000,0.000) circle (1.6pt);\node[font=\scriptsize\bfseries,text=navy] at ($(4.000,0.000)+(-73.6:0.32cm)$) {B};\filldraw[fill=white,draw=navy] (7.000,3.000) circle (1.6pt);\node[font=\scriptsize\bfseries,text=navy] at ($(7.000,3.000)+(-5.7:0.32cm)$) {C};\filldraw[fill=white,draw=navy] (4.000,6.000) circle (1.6pt);\node[font=\scriptsize\bfseries,text=navy] at ($(4.000,6.000)+(69.0:0.32cm)$) {D};\filldraw[fill=white,draw=navy] (0.000,6.000) circle (1.6pt);\node[font=\scriptsize\bfseries,text=navy] at ($(0.000,6.000)+(139.1:0.32cm)$) {E};\end{tikzpicture}\end{center}\begin{tasks}(1)
  \task Determina el sistema d'inequacions que defineix la regió.
  \task Calcula el màxim i el mínim de $f(x,y)=2x+y$ en la regió i on s'assoleixen.
  \task Estudia en quins punts de la regió assoleix el màxim i el mínim la funció $h(x,y)=x+2y$. Alguna de les dues solucions és única?
  \task S'afegeix la restricció $x\leq5$. Dibuixa la nova regió, calcula els vèrtexs nous i troba el nou màxim de $f$.
\end{tasks}\recorda{Per a escriure la inequació d'un costat, calcula la recta que passa per dos vèrtexs i tria el sentit de la desigualtat amb un punt de dins de la regió.}


\newpage
\begin{tcolorbox}[colback=teal,colframe=teal,coltext=white,arc=2mm,boxrule=0pt,left=3mm,right=3mm,top=1.5mm,bottom=1.5mm]
{\Large\bfseries\faCheckCircle\ \ Solucions}
\end{tcolorbox}
\setcounter{exer}{0}

\seccio{1}{Regió factible, optimització i infinites solucions}

\ex{1} \textbf{a)} Amb el punt de prova $(0,0)$ es compleixen totes les inequacions: la regió és el costat de cada recta on queda l'origen.\[\left\{\begin{aligned}-x+y&=3\\y&=-1\end{aligned}\right.\Rightarrow A(-4,-1)\qquad\left\{\begin{aligned}y&=-1\\2x-y&=9\end{aligned}\right.\Rightarrow B(4,-1)\]\[\left\{\begin{aligned}2x-y&=9\\x+2y&=12\end{aligned}\right.\Rightarrow C(6,3)\qquad\left\{\begin{aligned}x+2y&=12\\-x+y&=3\end{aligned}\right.\Rightarrow D(2,5)\]\begin{center}\begin{tikzpicture}[x=0.4769cm,y=0.4769cm,baseline=(current bounding box.center)]\begin{scope}\clip (-5,-2) rectangle (8,7);\foreach \i in {-4,-2,...,8} \draw[muted!25,very thin] (\i,-2)--(\i,7);\foreach \i in {-2,0,...,7} \draw[muted!25,very thin] (-5,\i)--(8,\i);\fill[teal!22] (4.0000,-1.0000)--(6.0000,3.0000)--(2.0000,5.0000)--(-4.0000,-1.0000)--cycle;\end{scope}\draw[-latex,thick] (-5,0)--(8.5,0) node[below left=-1pt,font=\scriptsize]{$x$};\draw[-latex,thick] (0,-2)--(0,7.5) node[below left=-1pt,font=\scriptsize]{$y$};\node[below,font=\tiny,text=muted] at (-4,0) {$-4$};\node[below,font=\tiny,text=muted] at (-2,0) {$-2$};\node[below,font=\tiny,text=muted] at (2,0) {$2$};\node[below,font=\tiny,text=muted] at (4,0) {$4$};\node[below,font=\tiny,text=muted] at (6,0) {$6$};\node[below,font=\tiny,text=muted] at (8,0) {$8$};\node[left,font=\tiny,text=muted] at (0,-2) {$-2$};\node[left,font=\tiny,text=muted] at (0,2) {$2$};\node[left,font=\tiny,text=muted] at (0,4) {$4$};\node[left,font=\tiny,text=muted] at (0,6) {$6$};\begin{scope}\clip (-5,-2) rectangle (8,7);\draw[accent,line width=1.1pt] (-6.0000,-3.0000)--(5.0000,8.0000);\draw[blue!70!black,line width=1.1pt] (-6.0000,9.0000)--(9.0000,1.5000);\draw[violet,line width=1.1pt] (3.0000,-3.0000)--(8.5000,8.0000);\draw[orange!85!black,line width=1.1pt] (-6.0000,-1.0000)--(9.0000,-1.0000);\end{scope}\filldraw[fill=white,draw=navy] (4.0000,-1.0000) circle (1.5pt);\node[font=\tiny,inner sep=1.5pt,fill=white,fill opacity=.75,text opacity=1] at ($(4.0000,-1.0000)+(-51.3:0.36cm)$) {$(4,\,-1)$};\filldraw[fill=white,draw=navy] (6.0000,3.0000) circle (1.5pt);\node[font=\tiny,inner sep=1.5pt,fill=white,fill opacity=.75,text opacity=1] at ($(6.0000,3.0000)+(20.6:0.36cm)$) {$(6,\,3)$};\filldraw[fill=white,draw=navy] (2.0000,5.0000) circle (1.5pt);\node[font=\tiny,inner sep=1.5pt,fill=white,fill opacity=.75,text opacity=1] at ($(2.0000,5.0000)+(90.0:0.36cm)$) {$(2,\,5)$};\filldraw[fill=white,draw=navy] (-4.0000,-1.0000) circle (1.5pt);\node[font=\tiny,inner sep=1.5pt,fill=white,fill opacity=.75,text opacity=1] at ($(-4.0000,-1.0000)+(-157.4:0.36cm)$) {$(-4,\,-1)$};\end{tikzpicture}\end{center}\textbf{b)} $P(0,4)$: $-0+4=4\not\leq3$, \rest{no} és de la regió. \quad $Q(3,2)$: $-1\leq3$, $7\leq12$, $4\leq9$, $2\geq-1$: \rest{sí} que hi és.\par \textbf{c)} La regió és fitada; avaluem $f$ en els vèrtexs:\begin{center}\small\begin{tabular}{cc}\hline Vèrtex & $f=3x+2y$\\\hline $(-4,\,-1)$ & $-14$ \ $\leftarrow$\\$(4,\,-1)$ & $10$\\$(6,\,3)$ & $24$ \ $\leftarrow$\\$(2,\,5)$ & $16$\\\hline\end{tabular}\end{center}Màxim: $\res{f(6,3)=24}$. \quad Mínim: $\res{f(-4,-1)=-14}$.\par \textbf{d)} Les rectes de nivell $kx+2y=c$ tenen pendent $-\frac k2$: han de ser paral·leles a un costat.\begin{itemize}\setlength\itemsep{1pt}\item Costat $CD$ ($x+2y=12$, pendent $-\frac12$): $k=1$. $g(A)=-6$, $g(B)=2$, $g(C)=g(D)=12$: \res{\text{màxim } 12 \text{ en tot el segment } CD}.\item Costat $DA$ ($-x+y=3$, pendent $1$): $k=-2$. $g(A)=g(D)=6$, $g(B)=-10$, $g(C)=-6$: \res{\text{màxim } 6 \text{ en tot el segment } DA}.\item Costats $AB$ ($k=0$) i $BC$ ($k=-4$): donen el \emph{mínim}, no el màxim. No serveixen.\end{itemize}\begin{center}\begin{tikzpicture}[x=0.4154cm,y=0.4154cm,baseline=(current bounding box.center)]\begin{scope}\clip (-5,-2) rectangle (8,7);\foreach \i in {-4,-2,...,8} \draw[muted!25,very thin] (\i,-2)--(\i,7);\foreach \i in {-2,0,...,7} \draw[muted!25,very thin] (-5,\i)--(8,\i);\fill[teal!22] (4.0000,-1.0000)--(6.0000,3.0000)--(2.0000,5.0000)--(-4.0000,-1.0000)--cycle;\end{scope}\draw[-latex,thick] (-5,0)--(8.5,0) node[below left=-1pt,font=\scriptsize]{$x$};\draw[-latex,thick] (0,-2)--(0,7.5) node[below left=-1pt,font=\scriptsize]{$y$};\node[below,font=\tiny,text=muted] at (-4,0) {$-4$};\node[below,font=\tiny,text=muted] at (-2,0) {$-2$};\node[below,font=\tiny,text=muted] at (2,0) {$2$};\node[below,font=\tiny,text=muted] at (4,0) {$4$};\node[below,font=\tiny,text=muted] at (6,0) {$6$};\node[below,font=\tiny,text=muted] at (8,0) {$8$};\node[left,font=\tiny,text=muted] at (0,-2) {$-2$};\node[left,font=\tiny,text=muted] at (0,2) {$2$};\node[left,font=\tiny,text=muted] at (0,4) {$4$};\node[left,font=\tiny,text=muted] at (0,6) {$6$};\begin{scope}\clip (-5,-2) rectangle (8,7);\draw[accent,line width=1.1pt] (-6.0000,-3.0000)--(5.0000,8.0000);\draw[blue!70!black,line width=1.1pt] (-6.0000,9.0000)--(9.0000,1.5000);\draw[violet,line width=1.1pt] (3.0000,-3.0000)--(8.5000,8.0000);\draw[orange!85!black,line width=1.1pt] (-6.0000,-1.0000)--(9.0000,-1.0000);\draw[navy,thick,densely dashed] (-6.0000,9.0000)--(9.0000,1.5000);\draw[accent,line width=3pt] (6.0000,3.0000)--(2.0000,5.0000);\end{scope}\filldraw[fill=white,draw=navy] (4.0000,-1.0000) circle (1.5pt);\node[font=\tiny,inner sep=1.5pt,fill=white,fill opacity=.75,text opacity=1] at ($(4.0000,-1.0000)+(-51.3:0.36cm)$) {$(4,\,-1)$};\fill[accent] (6.0000,3.0000) circle (2.3pt);\node[font=\tiny,inner sep=1.5pt,fill=white,fill opacity=.75,text opacity=1] at ($(6.0000,3.0000)+(20.6:0.36cm)$) {$(6,\,3)$};\fill[accent] (2.0000,5.0000) circle (2.3pt);\node[font=\tiny,inner sep=1.5pt,fill=white,fill opacity=.75,text opacity=1] at ($(2.0000,5.0000)+(90.0:0.36cm)$) {$(2,\,5)$};\filldraw[fill=white,draw=navy] (-4.0000,-1.0000) circle (1.5pt);\node[font=\tiny,inner sep=1.5pt,fill=white,fill opacity=.75,text opacity=1] at ($(-4.0000,-1.0000)+(-157.4:0.36cm)$) {$(-4,\,-1)$};\end{tikzpicture}\end{center}


\seccio{2}{Problema de maximitzar amb solucions enteres}

\ex{1} \textbf{a)} $x$ = caixes A, $y$ = caixes B. Maximitzar $B(x,y)=5x+4y$ subjecte a:\[\left\{\begin{aligned}2x+y&\leq100 \quad\text{(taronges)}\\x+2y&\leq80 \quad\text{(mandarines)}\\x+y&\geq20\\y&\geq10\\x&\geq0\end{aligned}\right.\]\textbf{b)} Vèrtexs: $A(0,20)$, $B(10,10)$, $C(45,10)$, $D(40,20)$ (tall de $2x+y=100$ i $x+2y=80$), $E(0,40)$.\begin{center}\begin{tikzpicture}[x=0.1033cm,y=0.1000cm,baseline=(current bounding box.center)]\begin{scope}\clip (0,0) rectangle (60,50);\foreach \i in {0,10,...,60} \draw[muted!25,very thin] (\i,0)--(\i,50);\foreach \i in {0,10,...,50} \draw[muted!25,very thin] (0,\i)--(60,\i);\fill[teal!22] (45.0000,10.0000)--(40.0000,20.0000)--(0.0000,40.0000)--(0.0000,20.0000)--(10.0000,10.0000)--cycle;\end{scope}\draw[-latex,thick] (0,0)--(62.5,0) node[below left=-1pt,font=\scriptsize]{$x$};\draw[-latex,thick] (0,0)--(0,52.5) node[below left=-1pt,font=\scriptsize]{$y$};\node[below,font=\tiny,text=muted] at (10,0) {$10$};\node[below,font=\tiny,text=muted] at (20,0) {$20$};\node[below,font=\tiny,text=muted] at (30,0) {$30$};\node[below,font=\tiny,text=muted] at (40,0) {$40$};\node[below,font=\tiny,text=muted] at (50,0) {$50$};\node[below,font=\tiny,text=muted] at (60,0) {$60$};\node[left,font=\tiny,text=muted] at (0,10) {$10$};\node[left,font=\tiny,text=muted] at (0,20) {$20$};\node[left,font=\tiny,text=muted] at (0,30) {$30$};\node[left,font=\tiny,text=muted] at (0,40) {$40$};\node[left,font=\tiny,text=muted] at (0,50) {$50$};\begin{scope}\clip (0,0) rectangle (60,50);\draw[accent,line width=1.1pt] (50.5000,-1.0000)--(24.5000,51.0000);\draw[blue!70!black,line width=1.1pt] (-1.0000,40.5000)--(61.0000,9.5000);\draw[violet,line width=1.1pt] (21.0000,-1.0000)--(-31.0000,51.0000);\draw[orange!85!black,line width=1.1pt] (-1.0000,10.0000)--(61.0000,10.0000);\end{scope}\filldraw[fill=white,draw=navy] (45.0000,10.0000) circle (1.5pt);\node[font=\tiny,inner sep=1.5pt,fill=white,fill opacity=.75,text opacity=1] at ($(45.0000,10.0000)+(-20.4:0.36cm)$) {$(45,\,10)$};\filldraw[fill=white,draw=navy] (40.0000,20.0000) circle (1.5pt);\node[font=\tiny,inner sep=1.5pt,fill=white,fill opacity=.75,text opacity=1] at ($(40.0000,20.0000)+(0.0:0.36cm)$) {$(40,\,20)$};\filldraw[fill=white,draw=navy] (0.0000,40.0000) circle (1.5pt);\node[font=\tiny,inner sep=1.5pt,fill=white,fill opacity=.75,text opacity=1] at ($(0.0000,40.0000)+(134.5:0.36cm)$) {$(0,\,40)$};\filldraw[fill=white,draw=navy] (0.0000,20.0000) circle (1.5pt);\node[font=\tiny,inner sep=1.5pt,fill=white,fill opacity=.75,text opacity=1] at ($(0.0000,20.0000)+(180.0:0.36cm)$) {$(0,\,20)$};\filldraw[fill=white,draw=navy] (10.0000,10.0000) circle (1.5pt);\node[font=\tiny,inner sep=1.5pt,fill=white,fill opacity=.75,text opacity=1] at ($(10.0000,10.0000)+(-132.9:0.36cm)$) {$(10,\,10)$};\end{tikzpicture}\end{center}\textbf{c)}\begin{center}\small\begin{tabular}{cc}\hline Vèrtex & $B=5x+4y$\\\hline $(0,\,20)$ & $80$\\$(10,\,10)$ & $90$\\$(45,\,10)$ & $265$\\$(40,\,20)$ & $280$ \ $\leftarrow$\\$(0,\,40)$ & $160$\\\hline\end{tabular}\end{center}\res{40 \text{ caixes A i } 20 \text{ caixes B}}, amb un benefici màxim de \res{280\ \text{€}}.\par \textbf{d)} Ara $B=6x+3y=3(2x+y)$, paral·lela a $2x+y=100$:\begin{center}\small\begin{tabular}{cc}\hline Vèrtex & $B=6x+3y$\\\hline $(0,\,20)$ & $60$\\$(10,\,10)$ & $90$\\$(45,\,10)$ & $300$ \ $\leftarrow$\\$(40,\,20)$ & $300$ \ $\leftarrow$\\$(0,\,40)$ & $120$\\\hline\end{tabular}\end{center}El màxim, 300~€, s'assoleix en tot el segment $CD$. Les solucions enteres són els punts de $2x+y=100$ amb $40\leq x\leq45$:\[(40,20),\ (41,18),\ (42,16),\ (43,14),\ (44,12),\ (45,10)\]Hi ha \res{6 \text{ maneres}} d'obtindre el benefici màxim de \res{300\ \text{€}}.


\seccio{3}{Problema de minimitzar en una regió no fitada}

\ex{1} \textbf{a)} $x$ = sacs de P1, $y$ = sacs de P2. Minimitzar $C(x,y)=40x+30y$ subjecte a:\[\left\{\begin{aligned}3x+y&\geq15 \quad\text{(proteïnes)}\\x+y&\geq9 \quad\text{(fibra)}\\x+3y&\geq15 \quad\text{(minerals)}\\x\geq0,\ y&\geq0\end{aligned}\right.\]\textbf{b)} Vèrtexs: $A(0,15)$, $B(3,6)$, $C(6,3)$, $D(15,0)$. La regió \rest{no és fitada}.\begin{center}\begin{tikzpicture}[x=0.3111cm,y=0.2706cm,baseline=(current bounding box.center)]\begin{scope}\clip (0,0) rectangle (18,17);\foreach \i in {0,5,...,18} \draw[muted!25,very thin] (\i,0)--(\i,17);\foreach \i in {0,2,...,17} \draw[muted!25,very thin] (0,\i)--(18,\i);\fill[teal!22] (15.0000,0.0000)--(18.0000,0.0000)--(18.0000,17.0000)--(0.0000,17.0000)--(0.0000,15.0000)--(3.0000,6.0000)--(6.0000,3.0000)--cycle;\end{scope}\draw[-latex,thick] (0,0)--(19.25,0) node[below left=-1pt,font=\scriptsize]{$x$};\draw[-latex,thick] (0,0)--(0,17.5) node[below left=-1pt,font=\scriptsize]{$y$};\node[below,font=\tiny,text=muted] at (5,0) {$5$};\node[below,font=\tiny,text=muted] at (10,0) {$10$};\node[below,font=\tiny,text=muted] at (15,0) {$15$};\node[left,font=\tiny,text=muted] at (0,2) {$2$};\node[left,font=\tiny,text=muted] at (0,4) {$4$};\node[left,font=\tiny,text=muted] at (0,6) {$6$};\node[left,font=\tiny,text=muted] at (0,8) {$8$};\node[left,font=\tiny,text=muted] at (0,10) {$10$};\node[left,font=\tiny,text=muted] at (0,12) {$12$};\node[left,font=\tiny,text=muted] at (0,14) {$14$};\node[left,font=\tiny,text=muted] at (0,16) {$16$};\begin{scope}\clip (0,0) rectangle (18,17);\draw[accent,line width=1.1pt] (5.3333,-1.0000)--(-1.0000,18.0000);\draw[blue!70!black,line width=1.1pt] (10.0000,-1.0000)--(-9.0000,18.0000);\draw[violet,line width=1.1pt] (-1.0000,5.3333)--(19.0000,-1.3333);\draw[navy,thick,densely dashed] (8.2500,-1.0000)--(-6.0000,18.0000);\end{scope}\filldraw[fill=white,draw=navy] (0.0000,15.0000) circle (1.5pt);\node[font=\tiny,inner sep=1.5pt,fill=white,fill opacity=.75,text opacity=1] at ($(0.0000,15.0000)+(127.5:0.36cm)$) {$(0,\,15)$};\fill[accent] (3.0000,6.0000) circle (2.3pt);\node[font=\tiny,inner sep=1.5pt,fill=white,fill opacity=.75,text opacity=1] at ($(3.0000,6.0000)+(180.0:0.36cm)$) {$(3,\,6)$};\filldraw[fill=white,draw=navy] (6.0000,3.0000) circle (1.5pt);\node[font=\tiny,inner sep=1.5pt,fill=white,fill opacity=.75,text opacity=1] at ($(6.0000,3.0000)+(-90.0:0.36cm)$) {$(6,\,3)$};\filldraw[fill=white,draw=navy] (15.0000,0.0000) circle (1.5pt);\node[font=\tiny,inner sep=1.5pt,fill=white,fill opacity=.75,text opacity=1] at ($(15.0000,0.0000)+(-30.1:0.36cm)$) {$(15,\,0)$};\end{tikzpicture}\end{center}\textbf{c)}\begin{center}\small\begin{tabular}{cc}\hline Vèrtex & $C=40x+30y$\\\hline $(0,\,15)$ & $450$\\$(3,\,6)$ & $300$ \ $\leftarrow$\\$(6,\,3)$ & $330$\\$(15,\,0)$ & $600$\\\hline\end{tabular}\end{center}Com que la funció creix en allunyar-se de l'origen, el mínim existeix i és en un vèrtex (la recta de nivell $40x+30y=300$, discontínua, deixa tota la regió al costat de cost més gran): \res{3 \text{ sacs de P1 i } 6 \text{ de P2}}, amb un cost mínim de \res{300\ \text{€}}.\par \textbf{d)} \rest{No hi ha cost màxim}: la regió no és fitada i $40x+30y$ creix indefinidament. Amb $(4,4)$: proteïnes $16\geq15$, fibra $8\not\geq9$, minerals $16\geq15$: \rest{no} cobriria la fibra.


\seccio{4}{Regió donada en un gràfic}

\ex{1} \textbf{a)} Vèrtexs: $A(0,2)$, $B(4,0)$, $C(7,3)$, $D(4,6)$, $E(0,6)$. Amb el punt interior $(2,3)$:\[AB:\ x+2y\geq4,\quad BC:\ x-y\leq4,\quad CD:\ x+y\leq10,\quad DE:\ y\leq6,\quad EA:\ x\geq0\]\textbf{b)} i \textbf{c)}\begin{center}\small\begin{tabular}{ccc}\hline Vèrtex & $f=2x+y$ & $h=x+2y$\\\hline $A(0,\,2)$ & $2$ & $4$\\$B(4,\,0)$ & $8$ & $4$\\$C(7,\,3)$ & $17$ & $13$\\$D(4,\,6)$ & $14$ & $16$\\$E(0,\,6)$ & $6$ & $12$\\\hline\end{tabular}\end{center}Màxim de $f$: $\res{17}$ en $C(7,3)$. Mínim de $f$: $\res{2}$ en $A(0,2)$.\par Màxim de $h$: $\res{16}$ només en $D(4,6)$ (solució única). Mínim de $h$: $\res{4}$ en $A$ i en $B$: com que $h$ és paral·lela al costat $x+2y=4$, el mínim s'assoleix en \rest{tots els punts del segment $AB$} (infinites solucions).\begin{center}\begin{tikzpicture}[x=0.5200cm,y=0.5200cm,baseline=(current bounding box.center)]\begin{scope}\clip (-1,-1) rectangle (9,7);\foreach \i in {0,2,...,9} \draw[muted!25,very thin] (\i,-1)--(\i,7);\foreach \i in {0,2,...,7} \draw[muted!25,very thin] (-1,\i)--(9,\i);\fill[teal!22] (7.0000,3.0000)--(4.0000,6.0000)--(0.0000,6.0000)--(0.0000,2.0000)--(4.0000,0.0000)--cycle;\end{scope}\draw[-latex,thick] (-1,0)--(9.5,0) node[below left=-1pt,font=\scriptsize]{$x$};\draw[-latex,thick] (0,-1)--(0,7.5) node[below left=-1pt,font=\scriptsize]{$y$};\node[below,font=\tiny,text=muted] at (2,0) {$2$};\node[below,font=\tiny,text=muted] at (4,0) {$4$};\node[below,font=\tiny,text=muted] at (6,0) {$6$};\node[below,font=\tiny,text=muted] at (8,0) {$8$};\node[left,font=\tiny,text=muted] at (0,2) {$2$};\node[left,font=\tiny,text=muted] at (0,4) {$4$};\node[left,font=\tiny,text=muted] at (0,6) {$6$};\begin{scope}\clip (-1,-1) rectangle (9,7);\draw[accent,line width=1.1pt] (-2.0000,3.0000)--(10.0000,-3.0000);\draw[blue!70!black,line width=1.1pt] (-2.0000,6.0000)--(10.0000,6.0000);\draw[violet,line width=1.1pt] (12.0000,-2.0000)--(2.0000,8.0000);\draw[orange!85!black,line width=1.1pt] (2.0000,-2.0000)--(12.0000,8.0000);\draw[accent,line width=3pt] (0.0000,2.0000)--(4.0000,0.0000);\end{scope}\fill[accent] (4.0000,6.0000) circle (2.3pt);\node[font=\tiny,inner sep=1.5pt,fill=white,fill opacity=.75,text opacity=1] at ($(4.0000,6.0000)+(69.0:0.36cm)$) {$(4,\,6)$};\filldraw[fill=white,draw=navy] (0.0000,6.0000) circle (1.5pt);\node[font=\tiny,inner sep=1.5pt,fill=white,fill opacity=.75,text opacity=1] at ($(0.0000,6.0000)+(139.1:0.36cm)$) {$(0,\,6)$};\filldraw[fill=white,draw=navy] (0.0000,2.0000) circle (1.5pt);\node[font=\tiny,inner sep=1.5pt,fill=white,fill opacity=.75,text opacity=1] at ($(0.0000,2.0000)+(-155.0:0.36cm)$) {$(0,\,2)$};\filldraw[fill=white,draw=navy] (4.0000,0.0000) circle (1.5pt);\node[font=\tiny,inner sep=1.5pt,fill=white,fill opacity=.75,text opacity=1] at ($(4.0000,0.0000)+(-73.6:0.36cm)$) {$(4,\,0)$};\filldraw[fill=white,draw=navy] (7.0000,3.0000) circle (1.5pt);\node[font=\tiny,inner sep=1.5pt,fill=white,fill opacity=.75,text opacity=1] at ($(7.0000,3.0000)+(-5.7:0.36cm)$) {$(7,\,3)$};\end{tikzpicture}\end{center}\textbf{d)} La recta $x=5$ retalla la punta $C$. Vèrtexs nous: $x=5$ amb $x-y=4\Rightarrow F(5,1)$; \ $x=5$ amb $x+y=10\Rightarrow G(5,5)$.\begin{center}\begin{tikzpicture}[x=0.5200cm,y=0.5200cm,baseline=(current bounding box.center)]\begin{scope}\clip (-1,-1) rectangle (9,7);\foreach \i in {0,2,...,9} \draw[muted!25,very thin] (\i,-1)--(\i,7);\foreach \i in {0,2,...,7} \draw[muted!25,very thin] (-1,\i)--(9,\i);\fill[teal!22] (5.0000,5.0000)--(4.0000,6.0000)--(0.0000,6.0000)--(0.0000,2.0000)--(4.0000,0.0000)--(5.0000,1.0000)--cycle;\end{scope}\draw[-latex,thick] (-1,0)--(9.5,0) node[below left=-1pt,font=\scriptsize]{$x$};\draw[-latex,thick] (0,-1)--(0,7.5) node[below left=-1pt,font=\scriptsize]{$y$};\node[below,font=\tiny,text=muted] at (2,0) {$2$};\node[below,font=\tiny,text=muted] at (4,0) {$4$};\node[below,font=\tiny,text=muted] at (6,0) {$6$};\node[below,font=\tiny,text=muted] at (8,0) {$8$};\node[left,font=\tiny,text=muted] at (0,2) {$2$};\node[left,font=\tiny,text=muted] at (0,4) {$4$};\node[left,font=\tiny,text=muted] at (0,6) {$6$};\begin{scope}\clip (-1,-1) rectangle (9,7);\draw[accent,line width=1.1pt] (-2.0000,3.0000)--(10.0000,-3.0000);\draw[blue!70!black,line width=1.1pt] (-2.0000,6.0000)--(10.0000,6.0000);\draw[violet,line width=1.1pt] (12.0000,-2.0000)--(2.0000,8.0000);\draw[orange!85!black,line width=1.1pt] (2.0000,-2.0000)--(12.0000,8.0000);\draw[teal,line width=1.1pt] (5.0000,-2.0000)--(5.0000,8.0000);\end{scope}\filldraw[fill=white,draw=navy] (4.0000,6.0000) circle (1.5pt);\node[font=\tiny,inner sep=1.5pt,fill=white,fill opacity=.75,text opacity=1] at ($(4.0000,6.0000)+(69.4:0.36cm)$) {$(4,\,6)$};\filldraw[fill=white,draw=navy] (0.0000,6.0000) circle (1.5pt);\node[font=\tiny,inner sep=1.5pt,fill=white,fill opacity=.75,text opacity=1] at ($(0.0000,6.0000)+(138.4:0.36cm)$) {$(0,\,6)$};\filldraw[fill=white,draw=navy] (0.0000,2.0000) circle (1.5pt);\node[font=\tiny,inner sep=1.5pt,fill=white,fill opacity=.75,text opacity=1] at ($(0.0000,2.0000)+(-156.0:0.36cm)$) {$(0,\,2)$};\filldraw[fill=white,draw=navy] (4.0000,0.0000) circle (1.5pt);\node[font=\tiny,inner sep=1.5pt,fill=white,fill opacity=.75,text opacity=1] at ($(4.0000,0.0000)+(-73.3:0.36cm)$) {$(4,\,0)$};\filldraw[fill=white,draw=navy] (5.0000,1.0000) circle (1.5pt);\node[font=\tiny,inner sep=1.5pt,fill=white,fill opacity=.75,text opacity=1] at ($(5.0000,1.0000)+(-49.4:0.36cm)$) {$(5,\,1)$};\fill[accent] (5.0000,5.0000) circle (2.3pt);\node[font=\tiny,inner sep=1.5pt,fill=white,fill opacity=.75,text opacity=1] at ($(5.0000,5.0000)+(39.8:0.36cm)$) {$(5,\,5)$};\end{tikzpicture}\end{center}\begin{center}\small\begin{tabular}{cc}\hline Vèrtex & $f=2x+y$\\\hline $(0,\,2)$ & $2$\\$(4,\,0)$ & $8$\\$(5,\,1)$ & $11$\\$(5,\,5)$ & $15$ \ $\leftarrow$\\$(4,\,6)$ & $14$\\$(0,\,6)$ & $6$\\\hline\end{tabular}\end{center}El nou màxim de $f$ és $\res{15}$, en $G(5,5)$.


\end{document}
